\setcounter{framenumber}{0}
%25

\begin{frame}
  \titlepage
\end{frame}

\begin{frame}{Learning Goals}
By the end of this lecture, students should be able to:

\begin{itemize}
    \item state the naive definition of probability;
    \item identify when the naive definition is appropriate;
    \item compute probabilities by counting favorable outcomes;
    \item use complements to simplify probability calculations;
\end{itemize}
\end{frame}

\lecturepart{Naive Probability}

\begin{frame}{From Events to Probabilities}
In the previous lecture:

\[
\Sample = \{\text{all possible outcomes}\},
\qquad
A \subseteq \Sample.
\]

\vspace{0.25cm}

Now we want to assign a number to an event:

\[
\Prob(A) = \text{probability that } A \text{ occurs}.
\]

\vspace{0.25cm}

\begin{keybox}{Question}
If every outcome in a finite sample space is equally likely, how should we define $\Prob(A)$?
\end{keybox}
\end{frame}

\begin{frame}{Definition: Naive Probability}
\begin{keybox}{Naive Definition of Probability}
Let $A$ be an event for an experiment with a finite sample space $\Sample$.

If all outcomes in $\Sample$ are equally likely, define

\[
\Prob_{\mathrm{naive}}(A)
=
\frac{|A|}{|\Sample|}
=
\frac{\text{number of outcomes favorable to } A}
{\text{total number of outcomes in } \Sample}.
\]
\end{keybox}

\vspace{0.25cm}

Here $|A|$ denotes the number of elements in the set $A$ (aka $card(A)$).
\vspace{1mm}

Question: Why does this definition is suitable for equally-likely outcomes?

\end{frame}

\begin{frame}{Naive Probability - Properties}


Naive probability satisfies

\[
0 \leq \Prob(A) \leq 1, \ \Prob(S)=1, \ \Prob(\emptyset)=0,
\qquad \text{for every } A \subseteq \Sample.
\]

Why? First, if $A \subseteq \Sample$, then
\[
0 \leq |A| \leq |\Sample|,
\]
so
\[
0 \leq \frac{|A|}{|\Sample|} \leq 1.
\]

Also,
\[
\Prob(\Sample)=\frac{|\Sample|}{|\Sample|}=1,
\qquad
\Prob(\emptyset)=\frac{0}{|\Sample|}=0.
\]
\end{frame}

\begin{frame}{Additivity for Disjoint Events}
Suppose $A,B \subseteq \Sample$ ,

Then, since
\[
|A \cup B| = |A| + |B| - |A\cap B|.
\]

Therefore,
\[
\Prob(A \cup B)
=
\frac{|A \cup B|}{|\Sample|}
=
\frac{|A|+|B|-|A\cap B|}{|\Sample|}
=
\frac{|A|}{|\Sample|}
+
\frac{|B|}{|\Sample|}
-
\frac{|A\cap B|}{|\Sample|}.
\]

Thus
\[
{\Prob(A \cup B)=\Prob(A)+\Prob(B)} -\Prob(A\cap B).
\]
A speical case is when $A$ and $B$ are exclusive ($A\cap B=\emptyset$), when we have 
\[
{\Prob(A \cup B)=\Prob(A)+\Prob(B)}
\]


\begin{keybox}{Looking ahead}
Later, we will define probability more generally using properties like these as axioms.  
Naive probability will then become a special case of the general theory.
\end{keybox}
\end{frame}

\begin{frame}{What the Definition Means}
The naive definition says:

\[
\Prob(A)
=
\frac{\text{how many outcomes make } A \text{ happen}}
{\text{how many outcomes are possible}}.
\]

\vspace{0.3cm}

This requires two assumptions:

\begin{enumerate}
    \item The sample space $\Sample$ is finite (we can't divide by infinity).
    \item All outcomes in $\Sample$ are equally likely.
\end{enumerate}

\vspace{0.3cm}

\begin{keybox}{Important}
The naive definition is not wrong.  
It is just limited. It should only be used when its assumptions are justified.
\end{keybox}
\end{frame}

\begin{frame}{Pebble World Interpretation}
Suppose each pebble has the same mass.

\[
\Prob(A)
=
\frac{\text{number of pebbles in } A}
{\text{number of pebbles in } \Sample}.
\]

\begin{center}
\begin{tikzpicture}[scale=0.9]
    \draw[rounded corners, thick] (-3.4,-1.55) rectangle (3.4,1.55);
    \node at (3.05,1.25) {$\Sample$};

    \foreach \x/\y in {-2.6/0.8,-1.8/0.9,-1.0/0.65,-0.2/0.85,0.8/0.75,1.7/0.9,2.5/0.75,
                       -2.2/0.0,-1.2/-0.1,-0.3/0.05,0.7/-0.1,1.7/0.0,2.6/0.05,
                       -2.5/-0.85,-1.5/-0.75,-0.4/-0.9,0.6/-0.75,1.6/-0.85,2.5/-0.75}
        \fill[UWGray] (\x,\y) circle (2.5pt);

    \draw[UWPurple, thick] (-1.75,0.05) ellipse (1.25 and 1.05);
    \node[UWPurple] at (-2.8,1.15) {$A$};
\end{tikzpicture}
\end{center}

\[
\Prob(A) = \frac{|A|}{|\Sample|}.
\]
\end{frame}

\begin{frame}{Example: One Fair Die}
Roll one fair six-sided die.

\[
\Sample = \{1,2,3,4,5,6\}.
\]

Let
\[
A = \{\text{roll is even}\} = \{2,4,6\}.
\]

Then

\[
|\Sample| = 6,
\qquad
|A| = 3.
\]

Therefore

\[
\Prob(A)
=
\frac{|A|}{|\Sample|}
=
\frac{3}{6}
=
\frac{1}{2}.
\]
\end{frame}

\begin{frame}{Example: One Card}
Draw one card from a standard well-shuffled $52$-card deck.

Let

\[
A = \{\text{card is an ace}\}.
\]

There are $4$ aces and $52$ cards total, so

\[
\Prob(A)
=
\frac{|A|}{|\Sample|}
=
\frac{4}{52}
=
\frac{1}{13}.
\]

\vspace{0.25cm}

Let

\[
H = \{\text{card is a heart}\}.
\]

Then

\[
\Prob(H)
=
\frac{13}{52}
=
\frac{1}{4}.
\]
\end{frame}

\begin{frame}{Example: Red or Ace}
Draw one card from a standard $52$-card deck.

Let

\[
R = \{\text{card is red}\},
\qquad
A = \{\text{card is an ace}\}.
\]

We want

\[
\Prob(R \cup A).
\]

\vspace{0.2cm}

There are:
\[
26 \text{ red cards},
\qquad
4 \text{ aces},
\qquad
2 \text{ red aces}.
\]

So

\[
|R \cup A|
=
26 + 4 - 2
=
28.
\]

Therefore

\[
\Prob(R \cup A)
=
\frac{28}{52}
=
\frac{7}{13}.
\]
\end{frame}


\lecturepart{Complements}

\begin{frame}{Complement Rule Under the Naive Definition}
Let $A \subseteq \Sample$.

The complement of $A$ is

\[
A^c = \{s \in \Sample : s \notin A\}.
\]

Since every outcome is either in $A$ or not in $A$,

\[
|A^c| = |\Sample| - |A|.
\]

Therefore,

\[
\Prob(A^c)
=
\frac{|A^c|}{|\Sample|}
=
\frac{|\Sample|-|A|}{|\Sample|}
=
1 - \frac{|A|}{|\Sample|}.
\]

So

\[
{\Prob(A^c)=1-\Prob(A).}
\]
\end{frame}

\begin{frame}{Why Complements Are Useful}
Sometimes it is easier to count when an event does \textbf{not} happen.

\vspace{0.25cm}

Instead of computing $\Prob(A)$ directly, compute:

\[
\Prob(A) = 1 - \Prob(A^c).
\]

\vspace{0.3cm}

\begin{keybox}{Strategy}
When a probability asks for phrases like

\[
\text{``at least one''}, \qquad \text{``at least two''}, \qquad \text{``not all''},
\]

first ask whether the complement is easier.
\end{keybox}
\end{frame}

\begin{frame}{Example: At Least One Head}
Flip a fair coin $5$ times.

Let

\[
A = \{\text{at least one head occurs}\}.
\]

The complement is

\[
A^c = \{\text{no heads occur}\}.
\]

There is only one outcome with no heads: $TTTTT.$

The total number of outcomes is $2^5 = 32.$

Therefore

\[
\Prob(A^c)=\frac{1}{32},
\qquad
\Prob(A)=1-\frac{1}{32}=\frac{31}{32}.
\]
\end{frame}

\begin{frame}{Example: At Least One Six}
Roll a fair die $4$ times.

Let

\[
A = \{\text{at least one six occurs}\}.
\]

It is easier to count

\[
A^c = \{\text{no six occurs}\}.
\]

Total outcomes: $6^4$. Outcomes with no six: $5^4$.

Therefore,
\[
\Prob(A)
=
1-\Prob(A^c)
=
1-\frac{5^4}{6^4}.
\]
\end{frame}

\lecturepart{Common Mistakes}



\begin{frame}{Example: Two Fair Dice}
Roll two fair dice.

A correct sample space is

\[
\Sample = \{(i,j): i,j \in \{1,2,3,4,5,6\}\}.
\]

Here $(i,j)$ means:

\[
i = \text{result of die 1},
\qquad
j = \text{result of die 2}.
\]

Thus

\[
|\Sample| = 36.
\]

\vspace{0.25cm}

\begin{keybox}{Important}
The outcomes are ordered pairs.  
For example, $(5,6)$ and $(6,5)$ are different outcomes.
\end{keybox}
\end{frame}

\begin{frame}{Example: Sum 11 versus Sum 12}
For two fair dice:

\[
\Sample = \{(i,j): i,j \in \{1,\dots,6\}\}.
\]

Let

\[
A = \{\text{sum is } 11\},
\qquad
B = \{\text{sum is } 12\}.
\]

Then

\[
A = \{(5,6),(6,5)\},
\qquad
B = \{(6,6)\}.
\]

So

\[
\Prob(A)=\frac{2}{36}=\frac{1}{18},
\]

while $\Prob(B)=\frac{1}{36}.$

Thus a sum of $11$ is twice as likely as a sum of $12$.
\end{frame}


\begin{frame}{Mini-Exercises}
\begin{enumerate}
    \item A fair die is rolled. What is the probability of rolling a number at least $5$?

    \item A fair coin is flipped $4$ times. What is the probability of getting all tails?

    \item A card is drawn from a standard deck. What is the probability that it is not a heart?

    \item Two fair dice are rolled. What is the probability that the sum is $7$?

    \item Why is the statement ``the probability of rain tomorrow is $1/2$ because either it rains or it does not'' invalid?
\end{enumerate}

\vspace{0.2cm}

\begin{livebox}{Answers}
\[
\frac{2}{6}=\frac13,
\qquad
\frac{1}{16},
\qquad
1-\frac{13}{52}=\frac34,
\qquad
\frac{6}{36}=\frac16.
\]
\end{livebox}
\end{frame}

\begin{frame}{Summary}
\begin{itemize}
    \item The naive definition is
    \[
    \Prob(A)=\frac{|A|}{|\Sample|}.
    \]

    \item It applies when $\Sample$ is finite and all outcomes are equally likely.

    \item Naive probability satisfies:
    \[
    0 \leq \Prob(A) \leq 1,
    \qquad
    \Prob(\Sample)=1,
    \qquad
    \Prob(\emptyset)=0,
    \]
 If $A \cap B=\emptyset$, then
    \[
    \Prob(A \cup B)=\Prob(A)+\Prob(B).
    \]

    \item Complements are often useful:
    \[
    \Prob(A^c)=1-\Prob(A).
    \]
\end{itemize}

\begin{keybox}{Next lecture}
Counting methods: multiplication rule, sampling with replacement, and sampling without replacement.
\end{keybox}
\end{frame}